\begin{abstract}
Activation sparsity offers a route to efficient inference, but its value depends
on where zeros occur and how training accommodates them. We study these
choices through matched interventions in transformer architecture and
optimization: post-hoc clipping, feed-forward activation pressure, broader
gating, and pressure at the gated sites. A detailed 14M-parameter study and
selected 70M experiments compare validation loss with model-wide activation
sparsity opportunity, the fraction of counted multiplications with an exact-zero
operand. At 14M, expanding pressure from the feed-forward hidden state to four
sites increases loss at every tested threshold. Adding gates to attention
queries, keys, and values changes the response to pressure: at the largest
threshold, seven-site pressure adds 12.10 percentage points of opportunity
with a 0.127 loss increase, compared with 2.50 points and 0.378 for four-site
pressure. Trained interventions extend the evaluated uniform-clipping frontier
at both sizes, with benefits depending on the allowed loss increase.
Operation-level counts explain the contribution of attention sparsity and
the effect of architecture on cross-size comparisons. A correctness-checked
70M implementation study yields slower sparse execution in all tested settings.
These single-seed results identify site selection and optimization as coupled
design choices, and motivate interventions that also control the structure
required for efficient execution.
\end{abstract}
